\section{The release encoder}\label{sec:method}

\paragraph{Setting.}
The input is a table $D_{\mathrm{fit}}$ of $n$ rows with $p$ feature columns and a declared target $y$; the task is regression, binary or multiclass classification with $K$ classes.
The encoder writes one artifact $a$ with $|a|\le B$, where $|a|$ counts every byte of the file and $B=\VThree{BudgetBytes}$ by default.
The sampler reads $a$ and a seed and emits rows $(\tilde x,\tilde y)$ in the raw column space.
Every statistic in $a$ is computed from $D_{\mathrm{fit}}$ alone; nothing is fitted on evaluation rows.

\subsection{Container and the two candidates}\label{sec:container}
Every \cli{auto-v8} artifact is a \texttt{DPA1} container: a \VThree{DpaHeaderBytes}-byte header holding a magic, an encoder tag and a weights fingerprint (all zero when no learned weights are needed), followed by the payload.
The payload is one of two candidates fitted on all of $D_{\mathrm{fit}}$:
\begin{itemize}
  \item \textbf{Format~1}, a compressed program~\citep{dope2026v1}: a safetensors file holding one compressed stream with quantile marginals on a normal-score grid, a regime mixture of Gaussian copulas, and a target model made of a residual network and boosted oblivious trees, with bit depths chosen by a rate--distortion search. It is fitted within $B-\VThree{DpaHeaderBytes}$ bytes.
  \item \textbf{\texttt{DSE8}}, preceded by a compact column-schema section (types and level names needed to write raw rows) and fitted within the room that remains. If the schema leaves fewer than \VThree{MinPacketRoom}~bytes, \cli{auto-v8} writes Format~1 directly.
\end{itemize}

\subsection{The \texttt{DSE8} packet and its byte accounting}\label{sec:codec}
A \texttt{DSE8} packet is a \VThree{DseHeaderBytes}-byte header, a sequence of sections in ascending tag order, each framed by a \VThree{TlvBytes}-byte tag-and-length prefix, and a \VThree{CrcBytes}-byte CRC32 over all preceding bytes (Table~\ref{tab:bytes}; header fields in Appendix~\ref{app:codec}).
For a file written by \cli{auto-v8} with \texttt{DSE8} chosen,
\begin{equation}\label{eq:bytes}
  |a| \;=\; D + s + H + C + \sum_{\text{sections } j} \bigl(T + \ell_j\bigr) \;\le\; B ,
\end{equation}
where $D$, $H$, $T$ and $C$ are the fixed sizes of Table~\ref{tab:bytes}, $s$ is the schema section and $\ell_j$ the payload of section $j$.

\begin{table}[t]
  \centering\small
  \caption{Byte accounting of a \texttt{DSE8} file. Fixed sizes and the planner's byte model are read from the Rust sources by the paper renderer; $p$ active columns ($p\le\VThree{MaxCols}$), copula rank $r\le\VThree{MaxRank}$, $\kappa\le\VThree{KnotCount}$ quantile knots, $o=K$ outputs for multiclass targets and $o=$ one output otherwise.}
  \label{tab:bytes}
  \begin{tabular}{@{}lll@{}}
    \toprule
    Component & Bytes & Content \\
    \midrule
    \texttt{DPA1} header ($D$) & \VThree{DpaHeaderBytes} & magic, encoder tag, weights fingerprint \\
    Column schema ($s$) & variable & column types and level names (shorter with \cli{-{}-anonymize}) \\
    \texttt{DSE8} header ($H$) & \VThree{DseHeaderBytes} & task, $K$, $p$, flags, $r$, $\kappa$, quantized scalars \\
    Section frame ($T$), per section & \VThree{TlvBytes} & tag and payload length \\
    \midrule
    \emph{Mandatory sections} & & \\
    Scale card & variable & target location and scale or class ids; per column a kind, centre and scale \\
    Marginals & $p\,(\VThree{MargBaseBytes}+\kappa)$ + typed tails & kind, missing rate, range, $\kappa$ knots; lattice or atom tails \\
    Copula loadings & $p\,r$ & one signed byte per loading \\
    Target head & $\VThree{LinFixedBytes} + o\,(\VThree{LinPerOutBytes}+p)$ & temperature; per output a scale, an intercept, $p$ signed weights \\
    Nominal groups & variable & one-hot groups with level probabilities, when present \\
    \midrule
    \emph{Optional sections, in this order} & & \\
    Regimes & variable & mixture weights and factor means \\
    Main effects & \VThree{EffectBytes} each, at most \VThree{MaxEffects} & column, output, scale, \VThree{EffectBins} binned values \\
    Pair interactions & \VThree{PairBytes} each, at most \VThree{MaxPairs} & column pair, output, product and XOR weights \\
    Oblivious trees & \VThree{TreeBytes} each, at most \VThree{MaxTrees} & depth \VThree{TreeDepth}: split columns and thresholds, scale, \VThree{TreeLeaves} leaves \\
    Residual model & variable & regression only: heteroscedastic weights, quantile shifts \\
    \midrule
    CRC32 trailer ($C$) & \VThree{CrcBytes} & over every preceding byte \\
    \bottomrule
  \end{tabular}
\end{table}

\paragraph{Planning and packing.}
Before fitting, a planner picks the active columns $p$, the knot count $\kappa$ and the rank $r$ so that the mandatory sections, estimated with the byte model of Table~\ref{tab:bytes}, leave room for optional sections; columns beyond the $p$ active ones are not modelled, and the sampler writes their stored median.
After fitting, the packer writes the mandatory sections and then each optional section in the fixed order of Table~\ref{tab:bytes}.
Effects, pairs and trees are ordered by fitted utility, and the packer keeps the longest prefix of each list that still fits, so a packet can be truncated to a smaller budget by dropping trailing items without refitting.
If the mandatory sections alone exceed the room, the fit is retried with a smaller planning budget; a candidate that never fits is not produced.

\paragraph{Decoding.}
Row $i$ of seed $\sigma$ is decoded from counter-based random streams keyed by $(\sigma, i, \text{stream}, \text{coordinate})$, so rows are independent, reproducible and can be generated in any order or in parallel.
The generative model is a factor copula:
\begin{align}
  f &\sim \mathcal N(0, I_r)\ \text{(optionally shifted by a sampled regime)}, &
  z_j &= L_j^\top f + \sigma_j \varepsilon_j, &
  x_j &= Q_j\bigl(\Phi(z_j)\bigr),
\end{align}
with loadings $L_j$, noise scales $\sigma_j$ and monotone quantile functions $Q_j$ given by the marginal knots and typed tails (lattice steps, atoms, zero inflation, missingness).
One-hot nominal groups are decoded jointly from their level probabilities, tilted by the latent scores: in each decoded row at most one indicator of a group is active, with all zero when the group's omitted baseline level is drawn.
The target is drawn last from logits
$\eta = b + W x + \sum_e g_e(z) + \sum_q h_q(z) + \sum_t \mathrm{tree}_t(z)$
of the linear head, the main effects, the pair interactions and the oblivious trees: a softmax or logistic draw at a fitted temperature for classification, and for regression the head's mean plus noise whose scale may depend on $z$ and whose quantiles are shifted by the residual section.

\subsection{Analytic fit from the fitting rows}\label{sec:fit}
The released encoder fits \texttt{DSE8} without learned weights.
\begin{enumerate}
  \item \textbf{Scale card and typed marginals.} Column kinds (continuous, lattice, discrete atoms, zero-inflated, constant) are inferred from $D_{\mathrm{fit}}$. Each column's encoded range runs between its empirical quantiles at levels $u$ and $1-u$, where $u=\min(\VThree{TailCap},\,k/\max(n_{\mathrm{obs}}+1,\VThree{TailMinDen}))$ for $n_{\mathrm{obs}}$ observed values and support floor $k=\VThree{SupportK}$; discrete values that occur at least $k$ times may instead be encoded as atoms, depending on the column's type.
  \item \textbf{Copula.} Rank-$r$ loadings are fitted to tie-aware normal scores, with loading norms capped below one so that every column keeps its own noise; a small mixture of factor means (regimes) is fitted when the budget allows.
  \item \textbf{Target head.} Ridge-penalized Newton (IRLS) fits the linear head, with the penalty chosen by out-of-fold cross-validation and the logit temperature calibrated out of fold.
  \item \textbf{Gated residual boosting.} Main effects, pairwise product and XOR interactions, and depth-\VThree{TreeDepth} oblivious trees are fitted in that order on the working residual. Because the decoder evaluates them on latent scores $z$, each real value is binned by the interval of $z$ that its tie group occupies, so fitting and decoding see the same bins. A component is kept only if it lowers a proper loss~\citep{gneiting2007strictly} on a deterministic inner validation split.
  \item \textbf{Intercept calibration.} After packing, the encoder decodes a fixed probe of synthetic feature rows from the packet itself and solves for intercepts so that the mean predicted class probabilities (for regression, the mean prediction) on those rows match the fitting rows. This corrects the shift between the real $x$ the head was fitted on and the synthetic $x$ it is applied to.
\end{enumerate}

\subsection{Per-table selection (\cli{auto-v8})}\label{sec:select}
Format~1 and \texttt{DSE8} are both fitted on all of $D_{\mathrm{fit}}$; a check on $D_{\mathrm{fit}}$ decides which one is written.
Rows are grouped into exact-duplicate groups over the feature and target columns.
If there are fewer than $G=\VThree{MinCheckGroups}$ groups, Format~1 is written (or \texttt{DSE8} if Format~1 does not fit).
Otherwise the groups are shuffled, ordered by target value for classification, and dealt alternately into halves $H_1, H_2$, so no duplicate group straddles the halves and class proportions are balanced.
On fold $f$, each candidate is refitted on one half, $\VThree{CheckSynthFactor}$ times as many rows as that half are sampled from it, and two internal auditors (a linear model and gradient-boosted trees) are trained on the sample and on the real half and scored on the other half.
The candidate's fold score $s^{(f)}$ is the minimum over auditors of the retention $(L_0-L_{\mathrm{syn}})/(L_0-L_{\mathrm{real}})$, with $L_0$ the null loss; if a classification fold has fewer than \VThree{MinTestPerClass} test rows in some class, the check is void and Format~1 is kept.
With fold means $\bar s$, \texttt{DSE8} is written if and only if
\begin{equation}\label{eq:select}
  \bar s_{\texttt{DSE8}} > \bar s_{\mathrm{F1}} + \delta
  \quad\text{and}\quad
  s^{(f)}_{\texttt{DSE8}} \ge s^{(f)}_{\mathrm{F1}} - \nu \ \text{ for both folds } f,
  \qquad \delta=\VThree{SelectionMargin},\ \nu=\VThree{FoldVeto}.
\end{equation}
The margin $\delta$ is the constant that \cli{-{}-format auto} also uses; the per-fold veto $\nu$ stops one favourable fold from carrying the decision.
When learned \texttt{DSE8} weights are supplied (\cli{-{}-v8-weights}), a third, neural candidate enters the same rule and must beat both others by $\delta$; the release ships no such weights, so by default \cli{auto-v8} chooses between Format~1 and the analytic packet only.
